\documentclass[a4paper,11pt,oneside,openany]{book}
\usepackage{styles/handbook}
\hypersetup{pdftitle={第十一批阶段样章：文本、序列与注意力}}
\begin{document}
\frontmatter
\hypersetup{pageanchor=false}
\begin{titlepage}
\vspace*{16mm}
{\sffamily\color{Teal}\large 机器学习与深度学习：图解速查手册\par}
\vspace{15mm}
{\sffamily\bfseries\color{Ink}\Huge 文本、序列与注意力\par}
\vspace{6mm}
{\sffamily\color{Ink}\LARGE 从词袋到Transformer\par}
\vspace{12mm}
\noindent\textcolor{Teal}{\rule{\textwidth}{1.2pt}}
\vspace{7mm}
\begin{keyidea}[第十一批阶段样章]
本批详细核查6张文本表示原图，并独立补写RNN、门控网络和Transformer。
用新图与算例连接Token、稀疏/稠密表示、共享状态、注意力与信息遮罩。
本文件是阶段成果；全书仍在连续修订。
\end{keyidea}
\vfill
原项目作者：张海飞、郭金珊、任露芹、李慧、朱靖、朱弘毅\par
\vspace{4mm}
2026年10月28日\par
\end{titlepage}
\hypersetup{pageanchor=true}
\tableofcontents
\mainmatter
\setcounter{chapter}{10}
\chapter{文本表示：从符号到模型输入}
\label{ch:text}

文本要先转成模型能处理的表示。词袋、TF-IDF与嵌入向量保留的信息不同，
并非一个必然比另一个更高级。先用简单基线理解任务，再决定是否需要序列与上下文模型。
本章先建立数值表示，下一章再讨论不同位置之间怎样交换信息。

\begin{keyidea}[符号编号没有大小距离含义]
Token ID是查表索引，不表示语义的高低远近。词典、分词规则与模型输入必须一致；
需要从当前数据学习的词典和文档频率，应在训练内拟合。
\end{keyidea}

\section{Tokenization、词干与停用词}
\label{sec:text-tokens}
Tokenization把文本切成处理单元（Token），可以是词、子词、字符等。
中文不必按空格切分，词干提取也不是所有语言和模型都需要的步骤。

\begin{figure}[htbp]
\centering\begin{tikzpicture}[node distance=8mm]
\node[box,text width=25mm] (text) {文本\\我爱猫};
\node[box,text width=30mm,right=of text] (tok) {Token序列\\我 / 爱 / 猫};
\node[box,text width=23mm,right=of tok] (ids) {ID序列\\1 / 2 / 3};
\node[box,text width=30mm,right=of ids] (emb) {嵌入查表\\$3\times d$};
\draw[flow] (text)--(tok);\draw[flow] (tok)--(ids);\draw[flow] (ids)--(emb);
\end{tikzpicture}

\caption{一个人工分词与词典例子，词典含PAD=0、我=1、爱=2、猫=3。实际系统可采用不同切分。}
\par\smallskip
\readfigure{三项ID查三行嵌入，得到三条向量。ID为3不表示“猫”是“我”的三倍。
验证与上线必须沿用同一词典及未知Token处理规则。}
\end{figure}

\begin{center}
\begin{tabularx}{\textwidth}{P{35mm}Y}
\toprule
处理 & 作用与限制\\
\midrule
词干提取 Stemming & 按规则缩减词形，输出未必是语言学词根或完整单词。\\
词形还原 Lemmatization & 试图还原词元，可能依赖词性或语言资源。\\
停用词 Stop words & 某些任务过滤高频功能词；否定、语气等也可能是关键信息，不能默认全删。\\
\bottomrule
\end{tabularx}
\end{center}

\quick{“不喜欢”中的“不”可能决定情感标签；文本清洗必须与任务语义一致。}

\section{词袋与TF-IDF：可解释的稀疏基线}
\label{sec:text-tfidf}
词袋（Bag of Words）把文档表示为词典上的计数或出现指示向量。
它忽略词序：“猫追狗”与“狗追猫”在相同分词下可以具有相同词袋，但含义不同。
指示值0/1不是一个“二元分类器”，只是特征数值。

\begin{center}
\begin{tabular}{lrrrr}
\toprule
文档/词典 & 猫 & 狗 & 看 & 跑\\
\midrule
D1：猫 看 猫 & 2 & 0 & 1 & 0\\
D2：狗 看 & 0 & 1 & 1 & 0\\
D3：猫 跑 & 1 & 0 & 0 & 1\\
文档频率$df$ & 2 & 1 & 2 & 1\\
\bottomrule
\end{tabular}
\end{center}

TF-IDF组合词频与逆文档频率。本例使用原始计数$tf$、自然对数和平滑式
\[
 idf(t)=\log\frac{1+N}{1+df(t)}+1,\qquad v_{d,t}=tf(t,d)\,idf(t),
\]
其中$N=3$，$df$按出现该词的\textbf{文档数}计，不是总次数。
“猫”和“看”的$idf=1+\log(4/3)\approx1.288$，另外两项为$1+\log2\approx1.693$。
若再作L2归一化，D1变为$[2/\sqrt5,0,1/\sqrt5,0]$。
输入上述计数矩阵时，平滑、原始计数与L2归一化口径对应常见\texttt{TfidfTransformer}设置；
这里不假定默认分词器会按单个汉字切分。
若一篇新文档没有任何词落在固定词典内，它会形成全零向量，不能直接除以0来归一化。

不同实现有不同TF、平滑与归一化约定，应报告完整设置。
常见词权重降低不代表它没有语义价值，稀有词权重高也不证明它具有预测或因果作用。

\quick{词典与IDF从训练数据学习，随后固定用于验证、测试和新文档。}

\section{嵌入与Word2Vec：学习稠密向量}
\label{sec:text-embedding}
嵌入（Embedding）用矩阵$E\in\R^{V\times d}$为$V$个Token存储$d$维向量。
长度$T$的ID序列查表得到$T\times d$矩阵，批量后常为$B\times T\times d$。
仅嵌入矩阵就有$Vd$个参数；PAD、未知词和是否冻结按实现设计。

Word2Vec通过局部上下文预测等目标学习词向量，典型形式包括CBOW与Skip-gram。
上下文相似可能产生向量相近，但相似性和类比关系是数据与训练的结果，
不是对任意词都成立的语义定理。词向量还可能保留语料中的偏差。

\begin{center}
\begin{tabularx}{\textwidth}{P{35mm}YY}
\toprule
表示 & 同一Token在不同句子中 & 常见考虑\\
\midrule
词袋/TF-IDF & 对应固定词典列。 & 简单、稀疏；通常忽略顺序。\\
静态词嵌入 & 查同一行向量。 & 稠密；多义词共享初始表示。\\
上下文表示 & 经序列/注意力层后可随上下文变化。 & 依结构处理次序与可见范围；训练和使用成本通常更高。\\
\bottomrule
\end{tabularx}
\end{center}

预训练嵌入可冻结后作为特征，也可以随任务微调。预处理应与预训练模型约定一致，
不能任意删词、改变Token划分后仍假定原模型语义不变。

\begingroup\small
\noindent\textbf{参考：}Manning等，\textit{Introduction to Information Retrieval}（2008）第6章；
Mikolov等（2013）的Word2Vec工作；scikit-learn Text feature extraction文档。
具体TF-IDF口径见\href{https://scikit-learn.org/stable/modules/feature_extraction.html#tfidf-term-weighting}{词频加权与归一化说明}。
本章表格与流程图独立设计。
\par\endgroup

\chapter{序列与注意力：共享状态与动态汇总}
\label{ch:sequence}

序列模型需要处理位置之间的依赖。RNN按顺序更新状态，注意力根据当前表示计算位置之间的权重。
本章沿着输入形状、计算路径与可见信息，建立RNN、门控结构和Transformer的基本联系。

\begin{keyidea}[能看到什么，是模型定义的一部分]
离线阅读整段文本与根据历史预测未来，允许的信息不同。
双向处理、填充遮罩与因果遮罩必须围绕任务设定，不能仅把它们当作程序选项。
\end{keyidea}

\section{RNN：同一组权重逐位置更新状态}
\label{sec:seq-rnn}
基本循环神经网络（Recurrent Neural Network, RNN）可写为
\[
 \vct h_t=\tanh(W_x\vct x_t+W_h\vct h_{t-1}+\vct b).
\]
若输入维度$d$、状态维度$h$，则$W_x\in\R^{h\times d}$、$W_h\in\R^{h\times h}$。
同一组参数在所有位置复用，序列加长不自动增加这组参数数量。

\begin{figure}[H]
\centering\begin{tikzpicture}[node distance=12mm]
\node[box,text width=17mm] (h0) {$h_0$};
\node[box,text width=24mm,right=of h0] (h1) {RNN单元\\$h_1$};
\node[box,text width=24mm,right=of h1] (h2) {RNN单元\\$h_2$};
\node[box,text width=24mm,right=of h2] (h3) {RNN单元\\$h_3$};
\draw[flow] (h0)--(h1);\draw[flow] (h1)--(h2);\draw[flow] (h2)--(h3);
\node[lab,below=10mm of h1] (x1) {$x_1$};\node[lab,below=10mm of h2] (x2) {$x_2$};\node[lab,below=10mm of h3] (x3) {$x_3$};
\draw[flow] (x1)--(h1);\draw[flow] (x2)--(h2);\draw[flow] (x3)--(h3);
\node[lab,above=10mm of h2] {三个位置使用同一组权重，不是三个独立模型};
\end{tikzpicture}

\caption{展开三个位置只是展示计算依赖，横向箭头传递状态。}
\end{figure}
\readfigure{标量例子$h_t=\tanh(x_t+0.5h_{t-1})$，$h_0=0$、输入$[1,0,-1]$，
得到约$[0.762,0.363,-0.674]$。同一个输入值出现在不同历史后，可以产生不同状态。}

反向传播通过时间（Backpropagation Through Time, BPTT）沿展开图计算，
共享参数在各位置的贡献需要相加。长路径上的Jacobian连乘可能导致梯度消失或爆炸；
截断BPTT限制回传长度，裁剪可帮助控制爆炸，但两者都不是保证学到所有长期依赖。

\quick{双向RNN可以利用前后文，因而不能直接用于当时不可获得未来输入的实时预测。}

\clearpage
\section{LSTM与GRU：门控制信息更新}
\label{sec:seq-gates}
长短期记忆网络（Long Short-Term Memory, LSTM）的经典形式用Sigmoid门与候选状态控制更新：
\[
 \vct c_t=\vct f_t\odot\vct c_{t-1}+\vct i_t\odot\widetilde{\vct c}_t,
 \qquad\vct h_t=\vct o_t\odot\tanh(\vct c_t).
\]
遗忘门$f$、输入门$i$、输出门$o$分别由当前输入和前一状态的仿射变换经Sigmoid得到；
候选$\widetilde c$常用tanh。$c$为记忆单元，$h$为输出状态，$\odot$表示逐元素乘法。

\begin{center}
\begin{tabularx}{\textwidth}{P{31mm}Y}
\toprule
结构 & 主要区别\\
\midrule
基本RNN & 直接由输入与旧状态生成新状态，路径反复经过同类非线性。\\
LSTM & 有$c$与$h$两种状态，门控制旧记忆保留、候选写入与输出。\\
GRU & 通常只有一个状态，用更新门与重置门组合旧状态和候选。\\
\bottomrule
\end{tabularx}
\end{center}

一种GRU约定为$h_t=z_t\odot h_{t-1}+(1-z_t)\odot\widetilde h_t$，
候选依赖输入与$r_t\odot h_{t-1}$。不同文献可对$z$采用互补命名，读实现时要对照公式。
LSTM的$i$与$f$一般独立，不要求两者之和为1。

标量LSTM例子$c_{t-1}=2,f=0.9,i=0.2,\widetilde c=-0.5,o=0.8$，
则$c_t=1.7$，$h_t=0.8\tanh1.7\approx0.748$。
门控提供更适合保留状态的路径，但并不保证永远记住所有历史。

\clearpage
\section{注意力：查询键，再汇总值}
\label{sec:seq-attention}
缩放点积注意力（Scaled Dot-Product Attention）为
\[
 \operatorname{Attn}(Q,K,V)=\softmax\!\left(\frac{QK^\top}{\sqrt{d_k}}+M\right)V.
\]
Softmax逐行作用。$Q$是查询（Query），$K$是键（Key），$V$是值（Value），
$M$指定允许关注的位置；键与值的行数必须对应。

\begin{figure}[H]
\centering\begin{tikzpicture}[node distance=8mm]
\node[box,text width=25mm] (q) {查询$Q$\\寻找匹配};
\node[box,text width=30mm,right=of q] (score) {$QK^\top/\sqrt{d_k}$\\与键比较};
\node[box,text width=28mm,right=of score] (weight) {遮罩+Softmax\\每行权重};
\node[box,text width=30mm,right=of weight] (value) {权重乘$V$\\汇总值向量};
\draw[flow] (q)--(score);\draw[flow] (score)--(weight);\draw[flow] (weight)--(value);
\end{tikzpicture}

\caption{键参与计算匹配分数，值提供被加权汇总的内容；两者职责不同。}
\end{figure}
\readfigure{一个查询$q=\log3$、$d_k=1$，两个键为$[0,1]$，得分为$[0,\log3]$。
权重为$[1/4,3/4]$；若值为$(2,0)$和$(0,4)$，输出为$(0.5,3)$。}

若$Q\in\R^{T_q\times d_k}$、$K\in\R^{T_k\times d_k}$、$V\in\R^{T_k\times d_v}$，
权重形状是$T_q\times T_k$，输出为$T_q\times d_v$。
自注意力（Self-attention）的Q/K/V来自同一序列的不同投影，不要求三个矩阵数值相同。
交叉注意力可让查询与键/值来自不同序列。

\quick{权重随输入变化，是一种计算关系；不能直接把权重大小当作因果解释。}

\clearpage
\section{因果与填充遮罩：在Softmax之前处理}
\label{sec:seq-mask}
因果自注意力只允许位置$i$关注$j\leq i$。
常令$M_{ij}=0$表示允许，$M_{ij}=-\infty$表示屏蔽，再计算Softmax。
这使被屏蔽位置的权重为0，而允许位置重新归一化。

\begin{figure}[H]
\centering\begin{tikzpicture}[x=17mm,y=11mm]
\foreach \j in {1,2,3}{\node[lab] at (\j,0) {键\j};}
\foreach \i in {1,2,3}{
 \node[lab,anchor=east] at (.3,-\i) {查询\i};
 \foreach \j in {1,2,3}{
  \ifnum\j>\i\def\cc{gray!12}\def\val{0}\else\def\cc{Pale}\ifnum\i=1\def\val{1}\else\ifnum\i=2\def\val{1/2}\else\def\val{1/3}\fi\fi\fi
  \node[draw=Teal!30,fill=\cc,minimum width=16mm,minimum height=10mm] at (\j,-\i) {$\val$};}}
\node[lab] at (2,-4) {右上三角被遮罩；每行允许的权重和为1};
\end{tikzpicture}

\caption{三个位置的未遮罩得分全为0时，加入因果遮罩后得到的注意力权重。}
\end{figure}
\readfigure{第一行只能看第1位置，第二行平均看前两个，第三行平均看三个。
先做Softmax再简单把未来位置设0而不重新归一化，会得到另一种运算。}

填充遮罩（Padding Mask）用于排除PAD等无效位置，和因果遮罩解决不同问题。
损失或序列池化也要排除相应填充位置；不要只屏蔽键就默认所有输出均有效。
避免整行都被屏蔽，使Softmax没有可用位置；实际布尔遮罩的True含义还需对照接口。

\begin{pitfall}[训练可并行，不表示能偷看未来]
训练时可以同时计算所有位置，但因果遮罩仍限制每个位置的可见信息。
逐步生成则按已有前缀扩展，两种计算方式可以实现同一个信息约束。
\end{pitfall}

\clearpage
\section{多头注意力：多组投影再拼接}
\label{sec:seq-heads}
多头注意力（Multi-Head Attention, MHA）并行计算多个注意力头：
\[\mathrm{head}_i=\operatorname{Attn}(XW_i^Q,XW_i^K,XW_i^V).\]
再拼接并乘输出投影$W_O$。头之间通常有不同的可学习投影。

\begin{figure}[H]
\centering\begin{tikzpicture}[node distance=8mm and 11mm]
\node[box,text width=24mm] (input) {$X$\\$3\times4$};
\node[box,text width=27mm,right=of input,yshift=11mm] (h1) {头1输出\\$3\times2$};
\node[box,text width=27mm,right=of input,yshift=-11mm] (h2) {头2输出\\$3\times2$};
\node[box,text width=25mm,right=of $(h1)!0.5!(h2)$] (join) {拼接\\$3\times4$};
\node[box,text width=25mm,right=of join] (out) {$W_O$投影\\$3\times4$};
\draw[flow] (input)--(h1);\draw[flow] (input)--(h2);\draw[flow] (h1)--(join);\draw[flow] (h2)--(join);\draw[flow] (join)--(out);
\end{tikzpicture}

\caption{省略批量轴，长度3、表示维度4、两个等宽头的形状例子。}
\end{figure}
\readfigure{每头的键和值维度为2，输出各为$3\times2$，拼接恢复$3\times4$。
拼接不是取平均；$W_O$进一步混合不同头的输出。}

在常见等宽设置中，$d_{\mathrm{model}}$能被头数整除。
若Q、K、V和输出均从$d$维映到$d$维，忽略偏置，四组投影参数共$4d^2$；
这不是整个Transformer的参数数，还未包括前馈层、归一化和嵌入。

\begin{pitfall}[每个头不保证学到一种固定语言功能]
可以观察到不同注意模式，但不能预先指定“某头一定负责语法、某头一定负责语义”。
头的作用依赖任务、数据和训练结果。
\end{pitfall}

\clearpage
\section{Transformer块：注意力、前馈与残差}
\label{sec:seq-transformer}
Transformer把注意力、逐位置前馈层、残差和归一化组合起来。
下面选择一种Pre-LN示意：
$Z=X+\mathrm{MHA}(\mathrm{LN}(X))$，
$Y=Z+\mathrm{FFN}(\mathrm{LN}(Z))$。

\begin{figure}[H]
\centering\begin{tikzpicture}[node distance=11mm and 12mm]
\node[box,text width=20mm] (x) {$X$};
\node[box,text width=44mm,right=of x] (attn) {多头注意力\\$\mathrm{MHA}(\mathrm{LN}(X))$};
\node[circle,draw=Teal,fill=Pale,minimum size=8mm,right=of attn] (a1) {$+$};
\draw[flow] (x)--(attn);\draw[flow] (attn)--(a1);
\draw[flow] (x.north)--++(0,10mm)-|(a1.north);
\node[box,text width=20mm,below=22mm of x] (z) {$Z$};
\node[box,text width=44mm,right=of z] (ffn) {逐位置前馈层\\$\mathrm{FFN}(\mathrm{LN}(Z))$};
\node[circle,draw=Teal,fill=Pale,minimum size=8mm,right=of ffn] (a2) {$+$};
\node[box,text width=20mm,right=of a2] (y) {$Y$};
\draw[flow] (a1.south)--++(0,-8mm)-|(z.north);
\draw[flow] (z)--(ffn);\draw[flow] (ffn)--(a2);\draw[flow] (a2)--(y);
\draw[flow] (z.south)--++(0,-9mm)-|(a2.south);
\end{tikzpicture}

\caption{省略Dropout与部分实现细节的Pre-LN块；不是所有Transformer变体的唯一顺序。}
\end{figure}
\readfigure{注意力在位置间汇总；FFN对每个位置分别应用同一组权重，
典型为$d\to d_{\mathrm{ff}}\to d$的两次仿射变换与中间非线性。
残差支路保持对应形状，LayerNorm常按每个位置的特征轴计算。}

普通自注意力若没有位置或其他次序结构，对位置排列具有等变性。
序列模型通常加入位置表示，如可学习位置向量、正弦编码或相对位置机制。
它们提供次序信息，不能用“不同Token的编号大小”代替。

\quick{Transformer不是只有一个Softmax：表示、位置、遮罩、注意力、前馈与训练目标共同决定行为。}

\clearpage
\section{编码器、解码器与下一Token目标}
\label{sec:seq-encoder-decoder}
编码器常对完整输入作双向表示；自回归解码器按前缀预测下一Token；
编码器—解码器结构还用交叉注意力读取源表示。

\begin{figure}[H]
\centering\begin{tikzpicture}[node distance=15mm and 15mm]
\node[box,text width=25mm] (src) {源序列};
\node[box,text width=31mm,right=of src] (enc) {Encoder\\编码源信息};
\node[box,text width=31mm,right=of enc] (memory) {源表示\\作为$K,V$};
\node[box,text width=25mm,below=of src] (tgt) {目标前缀};
\node[box,text width=31mm,right=of tgt] (dec) {Decoder\\因果自注意力\\交叉注意力};
\node[box,text width=31mm,right=of dec] (out) {下一个Token\\的logits};
\draw[flow] (src)--(enc);\draw[flow] (enc)--(memory);
\draw[flow] (tgt)--(dec);\draw[flow] (dec)--(out);
\draw[flow] (memory.south)--++(0,-6mm)-|(dec.north);
\end{tikzpicture}

\caption{条件生成示意：目标查询读取源序列的键和值，同时屏蔽目标未来位置。}
\end{figure}
\readfigure{交叉注意力中的键/值来自编码器源表示，查询来自解码器目标状态。
源长度与目标长度可以不同，不能按逐位置一一配对理解。}

自回归训练通常将输入和目标错开一个位置：
\begin{center}
\begin{tabular}{lccc}
\toprule
位置 & 1 & 2 & 3\\
\midrule
输入 & BOS & 我 & 爱\\
监督目标 & 我 & 爱 & 猫\\
\bottomrule
\end{tabular}
\end{center}
位置3在训练时只看到BOS、我、爱，预测猫；不能把猫同时放在该位置输入中让模型直接复制。
教师强制（Teacher Forcing）使用真实历史目标作为训练前缀，生成时则使用已生成前缀。

\begin{pitfall}[预训练目标不同，使用方式也不同]
掩码Token恢复与因果下一Token预测有不同可见信息和输出目标。
不能仅凭名称把任何编码器模型当作已经训练好的自回归生成器。
\end{pitfall}

\clearpage
\section{比较架构，同时检查时间与内存}
\label{sec:seq-compare}
没有架构在所有任务上都最好。局部结构、序列长度、可用数据与部署需求会影响选择。

\begin{center}
\begin{tabularx}{\textwidth}{P{23mm}P{50mm}Y}
\toprule
架构 & 信息交互 & 常见考虑\\
\midrule
CNN & 共享局部窗口，叠层扩大依赖范围。 & 网格与局部模式；步幅、边界和尺度。\\
RNN、LSTM、GRU & 按顺序更新共享状态。 & 有序计算与长路径；可逐步维护状态。\\
Transformer & 根据输入计算注意力交互。 & 位置、遮罩、长序列开销与任务目标。\\
\bottomrule
\end{tabularx}
\end{center}

显式注意力权重的元素数为$B\,h\,T_q\,T_k$。
$B=2,h=4,T_q=T_k=128$时，共131072个元素；float32仅这部分占0.5MiB。
长度都变为256时，该部分变为2MiB，即4倍。
这不是总训练显存；优化实现可能避免完整保存权重矩阵，其他激活和参数仍需计入。

\subsection{自测与回看}
共享参数是否意味着每个位置输出相同？\textbf{不意味着。}
LSTM输入门与遗忘门是否必须互补？\textbf{不必。}
多头输出是否简单取平均？\textbf{通常是拼接再投影。}
训练并行是否允许看到未来标签？\textbf{不允许。}

\noindent\textbf{参考：}Goodfellow等《Deep Learning》第10章；Hochreiter与Schmidhuber（1997）的LSTM；
Cho等（2014）的门控结构；Vaswani等，\textit{Attention Is All You Need}（2017）。
图、形状与数值例子独立构造。

\end{document}
